
\section{总体分析}

本文围绕微构体中导电介质填充的四个递进问题展开研究。针对问题一，建立单颗粒导通判定的几何模型，推导圆柱体表面到电极面的最短距离公式，对附件1中12个给定配置逐一判定导通状态并统计导通概率。针对问题二，构建基于并查集的多颗粒连通性判定算法，引入空间哈希网格加速近邻搜索，对不同填充率进行蒙特卡洛模拟以估计导通概率，同时分析附件2中49颗粒固定配置的连通状态。针对问题三，在问题二建立的蒙特卡洛模拟框架基础上，采用二分搜索策略在颗粒数空间中搜索使导通概率不低于90\%的最小填充率，并通过大样本验证确保结果可靠性。针对问题四，引入金属B球体并扩展连通性判定算法以处理双金属体系，通过网格搜索和成本分析确定满足导通概率约束的最低成本配方。

各问题之间呈递进关系。问题一建立了基本的几何模型和导通判定逻辑，为后续多颗粒模拟奠定基础。问题二将单颗粒分析扩展至多颗粒系统，引入蒙特卡洛方法和连通性算法。问题三在问题二的模拟框架基础上增加了优化搜索维度。问题四进一步拓展至双金属体系，引入经济性约束。四个问题从简单到复杂、从分析到优化，共同构成了导电复合材料仿真优化的完整研究链条。
